\chapter{Tongue-Tie Division}

\section*{What Is This Surgery?}
Tongue-tie division, also known as frenotomy or frenuloplasty, is a surgical procedure aimed at correcting ankyloglossia, a condition characterized by an abnormally short or tight lingual frenulum that restricts the movement of the tongue. This band of tissue connects the underside of the tongue to the floor of the mouth, and when it is too tight, it can hinder essential functions such as breastfeeding, swallowing, and speech articulation. The surgery involves the precise incision or excision of the frenulum to restore normal tongue mobility, thereby alleviating associated functional difficulties. This procedure is particularly significant in infants and young children, as it can greatly improve their ability to latch during breastfeeding, enhance their feeding efficiency, and facilitate proper speech development. Typically performed in an outpatient setting, tongue-tie division is minimally invasive, has a high success rate, and is associated with low morbidity.

\section*{Alternative Names}
\begin{itemize}
  \item Frenotomy
  \item Frenulotomy
  \item Lingual Frenulum Release
  \item Frenuloplasty
  \item Tongue-Tie Release Surgery
  \item Ankyloglossia Release
\end{itemize}

\section*{Relevant Surgical Anatomy}
The surgical anatomy relevant to tongue-tie division includes the lingual frenulum, the submandibular glands, and the surrounding vascular and neural structures. The lingual frenulum is a thin band of tissue that extends from the floor of the mouth to the ventral surface of the tongue. Its anatomical variations can significantly affect tongue mobility. 

The arterial supply to the area is primarily from the lingual artery, a branch of the external carotid artery, which provides blood to the tongue and its associated structures. Venous drainage occurs via the lingual vein, which drains into the internal jugular vein. 

Nerve supply is provided by the hypoglossal nerve (CN XII), which innervates the majority of the tongue muscles, and the lingual nerve, a branch of the mandibular division of the trigeminal nerve (CN V3), which provides sensory innervation to the anterior two-thirds of the tongue. 

Lymphatic drainage from the tongue is primarily to the submandibular and cervical lymph nodes. Key surgical landmarks include:

\begin{itemize}
  \item **Lingual Frenulum:** The tissue band being incised.
  \item **Wharton's Ducts:** The submandibular ducts that open into the floor of the mouth, located laterally to the frenulum.
  \item **Sublingual Glands:** Located beneath the mucosa of the floor of the mouth, these glands can be affected during surgery if not carefully managed.
\end{itemize}

Understanding these anatomical relationships is crucial for minimizing complications during the procedure.

\section*{Surgical Principle \& Rationale}
The primary surgical principle of tongue-tie division is to release the restrictive lingual frenulum that limits the tongue's mobility. In cases of ankyloglossia, the frenulum may be too short, thick, or attached too close to the tip of the tongue, which prevents the tongue from moving freely. The rationale behind the procedure is to create a longer and more flexible connection between the tongue and the floor of the mouth, thereby restoring optimal tongue function.

The technical goals of the surgery include:

\begin{itemize}
  \item **Complete Release:** Ensuring that the frenulum is adequately incised to allow full range of motion of the tongue.
  \item **Preservation of Adjacent Structures:** Avoiding damage to critical structures such as the submandibular ducts and major blood vessels.
  \item **Functional Restoration:** Enabling improved breastfeeding, swallowing, and speech articulation by enhancing the tongue's mobility.
\end{itemize}

The procedure is designed to provide immediate functional benefits, allowing for better feeding dynamics in infants and improved speech capabilities in older children.

\section*{History \& Surgical Evolution}
The practice of dividing the lingual frenulum to treat tongue-tie has historical roots that date back centuries. Early records indicate that midwives used rudimentary tools to perform this intervention, often without formal medical training. In the 17th century, Dr. Alfred Velpeau documented the procedure in medical literature, emphasizing the importance of technique and the need for careful assessment of the frenulum's anatomy.

The 20th century saw significant advancements in the understanding and approach to frenotomy. The introduction of local anesthesia transformed the procedure, allowing for a more comfortable experience for patients. The advent of laser technology in recent years has further refined the technique, making it quicker and less traumatic. Contemporary practice now emphasizes a functional assessment of ankyloglossia, focusing on symptomatic cases rather than purely anatomical criteria, thus guiding the decision for surgical intervention.

\section*{Clinical Indications}
Tongue-tie division is indicated in cases where ankyloglossia causes significant functional impairment. Common clinical indications include:

\begin{itemize}
  \item **Breastfeeding Difficulties:** Ineffective latch, poor milk transfer, or maternal nipple pain.
  \item **Poor Infant Weight Gain:** Insufficient caloric intake due to ineffective suckling.
  \item **Maternal Nipple Trauma:** Persistent pain or injury during breastfeeding.
  \item **Speech Articulation Problems:** Difficulty pronouncing specific phonemes in older children.
  \item **Oral Hygiene Issues:** Inability to effectively clean teeth, increasing the risk of dental problems.
  \item **Mechanical Issues:** Difficulty performing oral tasks such as licking or playing wind instruments.
  \item **Gastrointestinal Symptoms:** Excessive air intake during feeding leading to discomfort.
  \item **Dental Spacing Issues:** Potential for gaps between teeth due to frenulum attachment.
\end{itemize}

These indications highlight the importance of addressing ankyloglossia to improve quality of life and functional outcomes.

\section*{Clinical Positioning}
Tongue-tie division is primarily considered a first-line treatment for symptomatic ankyloglossia. When infants or children present with clear functional impairments related to a restrictive lingual frenulum, frenotomy is typically the recommended intervention following a thorough clinical assessment. 

While non-surgical approaches, such as lactation consultation or speech therapy, may be attempted first, they are often not sufficient to resolve the underlying anatomical issue. In cases where conservative measures fail, surgical intervention becomes the primary definitive treatment. For older children with speech difficulties, concurrent speech therapy may be initiated to optimize articulation post-frenotomy, but the surgical release directly addresses the anatomical constraint.

\section*{Surgical Technique \& Procedure}
The procedure for tongue-tie division is generally performed in an outpatient setting and can be accomplished without general anesthesia, particularly in infants. The steps involved in the surgical technique are as follows:

\begin{itemize}
  \item **Anesthesia and Positioning:** Local anesthetic cream or spray may be applied, or the procedure may be performed without anesthesia due to minimal nerve supply in the frenulum. The infant is swaddled and held securely, with the head stabilized.
  \item **Visualization:** The clinician uses a finger or a grooved director to lift the tongue and expose the lingual frenulum. Adequate lighting is essential for clear visualization.
  \item **Incision:** Using sterile scissors, a scalpel, or a CO$_{2}$ laser, the frenulum is carefully incised along the midline, ensuring that adjacent structures such as Wharton's ducts are preserved.
  \item **Assessment of Release:** After the incision, the tongue's range of motion is assessed to ensure adequate release. The tongue should be able to elevate and extend freely.
  \item **Hemostasis:** Minor bleeding is controlled with direct pressure or laser cauterization. Sutures are rarely needed for simple frenotomies.
  \item **Post-Procedure Care:** The infant is encouraged to breastfeed immediately after the procedure to provide comfort and initiate tongue movement.
\end{itemize}

This technique is designed to be quick and effective, minimizing discomfort and promoting rapid recovery.

\section*{Preoperative Workup \& Preparation}
Preparation for tongue-tie division is generally straightforward, particularly for infants. Key components of the preoperative workup include:

\begin{itemize}
  \item **Clinical Assessment:** A thorough examination by a pediatrician or lactation consultant to confirm the diagnosis of symptomatic ankyloglossia.
  \item **Timing of Feeding:** Scheduling the procedure just before a feeding to facilitate immediate breastfeeding post-operation.
  \item **Informed Consent:** Parents or guardians must provide informed consent after discussing the procedure's benefits, risks, and alternatives.
  \item **Medical History Review:** Assessing for any bleeding disorders, current medications, or allergies to local anesthetics.
  \item **NPO Status (Rare):** For cases requiring sedation, specific NPO guidelines will be provided.
  \item **Antibiotics (Rare):** Prophylactic antibiotics are generally not necessary for simple frenotomy.
\end{itemize}

These preparatory steps ensure that the procedure is performed safely and effectively.

\section*{Contraindications \& Precautions}
While tongue-tie division is a safe procedure, certain conditions may contraindicate its performance.

\textbf{Absolute Contraindications:}
\begin{itemize}
  \item **Active Bleeding Disorder:** Uncontrolled coagulopathy increases the risk of hemorrhage.
  \item **Acute Infection:** An active oral infection may delay the procedure.
  \item **Unstable Medical Condition:** Severe underlying medical conditions that pose risks during surgery.
\end{itemize}

\textbf{Relative Contraindications and Precautions:}
\begin{itemize}
  \item **True Posterior Tongue-Tie:** More complex cases may require extensive frenuloplasty and careful assessment.
  \item **Prematurity or Low Birth Weight:** Requires careful consideration of the infant's stability.
  \item **Anatomical Anomalies:** Other oral or craniofacial anomalies may complicate the procedure.
  \item **Medications:** Use of anticoagulants should be discussed with the clinician.
  \item **Unclear Diagnosis:** Further evaluation is warranted if functional impairment is not clearly attributable to ankyloglossia.
\end{itemize}

Awareness of these contraindications helps to mitigate risks associated with the procedure.

\section*{Complications \& Risks}
Tongue-tie division is generally low-risk, but potential complications exist, including:

\begin{itemize}
  \item **Bleeding:** Minor bleeding is common and usually self-limiting; significant hemorrhage is rare.
  \item **Infection:** Wound infection is uncommon due to rapid healing and good blood supply.
  \item **Anesthesia Complications:** Risks associated with local anesthesia are minimal but can include allergic reactions.
  \item **Reattachment (Re-tie):** The frenulum may heal back together, necessitating a repeat procedure.
  \item **Pain and Discomfort:** Temporary discomfort is expected and can be managed with over-the-counter pain relievers.
  \item **Damage to Adjacent Structures:**
    \begin{itemize}
      \item **Submandibular Ducts:** Injury can lead to salivary gland issues.
      \item **Blood Vessels/Nerves:** Rare but possible injuries to lingual arteries or nerves.
    \end{itemize}
  \item **Scarring:** Minimal scarring may occur but rarely impacts function.
  \item **Incomplete Release:** Insufficient division may not resolve functional issues.
  \item **No Improvement in Symptoms:** Some patients may not experience symptom resolution despite a technically successful procedure.
\end{itemize}

Understanding these risks is essential for informed consent and patient management.

\section*{Postoperative Recovery Timeline}
Postoperative recovery following tongue-tie division is generally swift and uncomplicated. 

In the immediate postoperative period, patients are monitored in the Post-Anesthesia Care Unit (PACU) for any signs of bleeding or complications. During the first 24-48 hours, parents are advised to observe the wound site for any unusual symptoms, such as excessive bleeding or signs of infection. 

In the first week, infants are encouraged to breastfeed frequently, which aids in wound healing and provides comfort. Activity restrictions are minimal, and parents are instructed on performing gentle stretching exercises to prevent reattachment of the frenulum. 

Long-term recovery typically spans several weeks, during which the wound heals completely. Most infants show significant improvements in feeding dynamics and weight gain within the first two weeks, while older children may require additional speech therapy to address any persisting articulation issues.

\section*{Surgical Findings \& Pathology}
During the procedure, the surgeon documents the appearance of the frenulum and assesses the degree of restriction. The pathology report is not typically generated for simple frenotomies, as the procedure focuses on the mechanical release of the frenulum rather than the removal of pathological tissue. However, any abnormal findings, such as signs of infection or unusual anatomical variations, may be noted for future reference.

\section*{Expected Postoperative Course}
A successful postoperative course following tongue-tie division is characterized by several key indicators. 

Patients typically experience a resolution of the symptoms that prompted the procedure, such as improved breastfeeding dynamics and reduced maternal nipple pain. Infants should demonstrate enhanced suckling patterns and consistent weight gain. For older children, improvements in speech articulation are expected, with many achieving clearer pronunciation of previously difficult sounds. 

The wound site generally heals without complications, showing minimal scarring, and there is no evidence of reattachment of the frenulum. Overall, a normal postoperative course leads to a freely mobile tongue with no limitations in feeding, speech, or oral hygiene.

\section*{Postoperative Care \& Follow-up}
Postoperative care is crucial to ensure successful healing and functional improvement. 

Follow-up visits are typically scheduled within 1-2 weeks to assess wound healing and confirm adequate release. For breastfeeding infants, continued follow-up with a lactation consultant is recommended to optimize latch and address any residual feeding challenges. 

Parents are instructed on specific oral exercises or stretches to perform multiple times daily for several weeks to prevent reattachment of the frenulum. For older children with speech difficulties, referral to a speech-language pathologist may be necessary to retrain tongue muscles. 

Regular dental check-ups are important to monitor oral hygiene and dental development. Parents should be advised to contact their doctor if they observe signs of infection, significant bleeding, or if feeding difficulties persist.

\section*{Epidemiology}
Ankyloglossia, or tongue-tie, is a relatively common congenital anomaly, with reported incidence rates ranging from 3\% to 10\% of live births. This variability is often attributed to differing diagnostic criteria and assessment methods. 

The condition is observed more frequently in males than females, with a male-to-female ratio of approximately 2:1 to 3:1. While many cases of ankyloglossia are asymptomatic and do not require intervention, a significant proportion present with functional difficulties, particularly breastfeeding problems in infancy. 

The frequency of frenotomy procedures has increased in recent years, reflecting greater awareness among healthcare providers and parents regarding its potential impact on feeding and development. Mortality directly attributable to frenotomy is virtually non-existent, and serious morbidity is exceedingly rare.

\section*{Outcomes \& Prognosis}
The outcomes following tongue-tie division are generally excellent, with high success rates in resolving the functional issues for which the procedure was performed. 

For infants with breastfeeding difficulties, studies report improvement in latch and reduction in maternal nipple pain in 80-95\% of cases. Infant weight gain typically improves as well. For older children, significant improvement in speech articulation is observed, although some may require adjunctive speech therapy to fully correct long-standing patterns. 

Recurrence of tongue-tie due to reattachment is a known risk but is minimized with diligent post-operative stretching exercises. Long-term prognosis is overwhelmingly positive, with most individuals experiencing no further issues related to their tongue mobility or function after a successful frenotomy.

\section*{References}
\begin{enumerate}
  \item MedlinePlus. Tongue-tie surgery. U.S. National Library of Medicine; 2024. Available from: \url{https://medlineplus.gov/tonguetie.html}
  
  \item Kliegman RM, St. Geme JW, Blum NJ, Shah SS, Tasker RC, Wilson KM. Nelson Textbook of Pediatrics. 22nd ed. Elsevier; 2024.
  
  \item American Academy of Otolaryngology--Head and Neck Surgery. Clinical Consensus Statement: Ankyloglossia in Children. Otolaryngology--Head and Neck Surgery. 2020;162(5):597-611.
  
  \item Lawrence RA, Lawrence RM. Breastfeeding: A Guide for the Medical Profession. 9th ed. Elsevier; 2021.
\end{enumerate}

\clearpage